\section{Arquitectura del Sistema}
\label{sec:arquitectura}

\subsection{Visió general}

El sistema Dal-i segueix una arquitectura de microserveis \textit{cloud-native} desplegada completament sobre Google Cloud Platform (GCP). La decisió de separar les operacions lleugeres (historial, càrrega d'imatges, parla i traducció) en Cloud Functions~\cite{cloudfunctions} independents, mantenint el pipeline de visió computacionalment intens en un servei Cloud Run~\cite{cloudrun} dedicat, permet escalar cada component de forma independent i reduir costos quan el sistema no està en ús.

La \cref{fig:arquitectura} mostra el diagrama d'arquitectura del sistema:

\begin{figure}[H]
\centering
\begin{tikzpicture}[
    node distance=0.8cm and 1.4cm,
    box/.style={draw, rounded corners=4pt, minimum width=3.8cm,
                minimum height=0.9cm, align=center, font=\small\sffamily,
                fill=dalblue!8, draw=dalblue!40},
    storage/.style={box, fill=gray!10, draw=gray!40},
    arr/.style={-{Stealth[length=5pt]}, thick, color=dalblue!70},
    lbl/.style={font=\scriptsize\sffamily\color{gray}, midway},
  ]

  \node[font=\scriptsize\sffamily\color{gray!80}, align=center] (auth)
    {\faLock\; Firebase Anonymous Auth};

  \node[box, below=0.3cm of auth] (fe)
    {\faDesktop\\[2pt]\textbf{Frontend}\\{\scriptsize Next.js · React · TS}};

  \node[box, below=1.0cm of fe] (cf)
    {\faCloud\\[2pt]\textbf{Cloud Functions}\\
     {\scriptsize history · upload · speech}};

  \node[box, right=2.2cm of cf] (cr)
    {\faCog\\[2pt]\textbf{Cloud Run}\\{\scriptsize FastAPI · Vision}};

  \node[storage, below=1.0cm of cf] (ai)
    {\faBrain\\[2pt]\textbf{Vertex AI · Gemini}\\
     {\scriptsize Speech · TTS · Translate}};

  \node[storage, below=1.0cm of cr] (fs)
    {\faDatabase\\[2pt]\textbf{Firestore + GCS}};

  \draw[arr] (fe)  -- node[lbl, left]{REST} (cf);
  \draw[arr] (cf)  -- node[lbl, above]{REST} (cr);
  \draw[arr] (cf)  -- (ai);
  \draw[arr] (cr)  -- (fs);
  \draw[arr] (cr.south) to[out=-120, in=30] (ai.east);

\end{tikzpicture}
\caption{Diagrama d'arquitectura del sistema Dal-i sobre Google Cloud Platform.}
\label{fig:arquitectura}
\end{figure}

\subsection{Frontend}

El frontend és una aplicació \textbf{Next.js 16}~\cite{nextjs} (React 19~\cite{react}, TypeScript, Tailwind CSS 4~\cite{tailwind}) compilada com a exportació estàtica (\texttt{output: "export"}). El bundle resultant es serveix directament des del contenidor Cloud Run, sense necessitat d'un servidor de renderització independent.

La comunicació amb el backend es centralitza al mòdul \texttt{app/lib/api.ts}, que s'encarrega d'injectar automàticament el token de Firebase Anonymous Auth~\cite{firebase} a totes les peticions com a capçalera \texttt{Authorization: Bearer <token>}. L'autenticació és completament transparent per a l'usuari: el SDK de Firebase crida \texttt{signInAnonymously()} en el primer muntatge del component i retorna un UID estable i persistent per al navegador.

\subsection{Cloud Run — Backend de Visió}

El servei Cloud Run~\cite{cloudrun} (\texttt{dal-i-api}, regió \texttt{europe-west1}) és un contenidor Docker que executa una API REST amb FastAPI~\cite{fastapi}. Conté el pipeline de processament d'imatges (computacionalment intens) i serveix els endpoints principals:

\begin{itemize}
  \item \texttt{GET /health}: comprovació d'estat del servei.
  \item \texttt{POST /process}: pipeline complet per a selecció d'estil per botons.
  \item \texttt{POST /process/voice}: pipeline per a entrada de veu (compatible amb Raspberry Pi).
  \item \texttt{POST /process/text}: pipeline per a entrada de text (compatible amb Raspberry Pi).
\end{itemize}

Els endpoints \texttt{/process} requereixen un token de Firebase vàlid (o retornen \texttt{uid="device"} per a clients sense autenticació, com la Raspberry Pi). Tots els endpoints estan protegits per un limitador de velocitat en memòria (\textit{sliding window}, 20 peticions/minut per IP).

El Dockerfile és multi-etapa: la primera etapa compila el frontend amb Node.js 20, i la segona copia el bundle estàtic i el codi Python al contenidor final basat en Python 3.12 slim.

\subsection{Cloud Functions}

Tres Cloud Functions~\cite{cloudfunctions} de 2a generació gestionen les operacions lleugeres:

\subsubsection*{history}
Gestiona el CRUD de l'historial de generacions per usuari sobre Firestore~\cite{firestore}. Suporta:
\begin{itemize}
  \item \texttt{GET}: llistat paginat (cursor basat en \texttt{start\_after(snapshot)}, 12 elements per pàgina).
  \item \texttt{DELETE ?id=\{docId\}}: esborra un element de Firestore \textbf{i} les imatges associades de Cloud Storage~\cite{gcs}.
  \item \texttt{DELETE ?deleteAll=true}: esborra tots els elements de l'usuari en lots de fins a 500 (límit de batch de Firestore) i els blobs de GCS en cascada.
\end{itemize}

\subsubsection*{upload}
Accepta una imatge per POST (\texttt{multipart/form-data}), la valida (màxim 10 MB, MIME \texttt{image/*}) i la puja al bucket \texttt{dal-i-bucket} de Cloud Storage~\cite{gcs} amb accés públic. Retorna l'URL pública i el nom del fitxer.

\subsubsection*{speech}
Agrupa tres operacions de parla i traducció en un sol endpoint amb paràmetre \texttt{?action=}:
\begin{itemize}
  \item \texttt{transcribe}: rep un fitxer d'àudio WebM/Opus, el transcriu amb Cloud Speech-to-Text~\cite{stt} en l'idioma seleccionat, tradueix la transcripció a l'anglès (Cloud Translation~\cite{translate}), i extreu l'estil artístic.
  \item \texttt{tts}: rep text en anglès i un codi d'idioma destí, tradueix el text a l'idioma de l'usuari, genera l'àudio MP3 amb Cloud TTS~\cite{tts} (veus Journey o WaveNet) i el retorna en base64.
  \item \texttt{translate}: traducció de text simple entre qualsevol parell d'idiomes suportats.
\end{itemize}

Totes les Cloud Functions verifiquen el token de Firebase amb
\code{firebase\_admin.auth().verify\_id\_token(token)}
i afegeixen les capçaleres CORS necessàries per permetre crides des del domini de Cloud Run.

\subsection{Emmagatzematge i Bases de Dades}

\textbf{Cloud Storage}~\cite{gcs} (\code{dal-i-bucket}) emmagatzema les imatges d'entrada (\code{images/}), les imatges transformades per Gemini i els resultats JSON amb les coordenades (\code{results/}). Tots els blobs es fan públics en el moment de la pujada.

\textbf{Firestore}~\cite{firestore} emmagatzema els registres de l'historial seguint la ruta \codepath{generation_history/{uid}/items/{docId}}, on \code{uid} és l'identificador anònim de Firebase. Cada registre inclou estil, URLs d'imatge, coordenades, descripció artística i marca de temps.

\subsection{Autenticació}

S'utilitza \textbf{Firebase Anonymous Authentication}~\cite{firebase}: quan un usuari accedeix a l'aplicació per primera vegada, el SDK de Firebase crea silenciosament un compte anònim i emmagatzema el token al localStorage del navegador. Aquest UID és estable entre sessions (mentre no s'esborri la memòria del navegador) i permet mantenir l'historial personal sense que l'usuari hagi de registrar-se.

El backend verifica el token a totes les peticions; per als clients de maquinari (Raspberry Pi) que no implementen Firebase Auth, el sistema admet peticions sense token i les associa a l'UID genèric \texttt{"device"}.
